Geografická distribuce odborové hustoty v~EU27 ukazuje na klastrování: severské země a~Belgie udržují hustotu nad \SIrange{50}{60}{\percent}, zatímco středo- a~jihovýchodní Evropa (včetně ČR) se pohybuje pod \SIrange{15}{20}{\percent}.
Srovnání s~mapou \ac{KS} pokrytí (obr.~\ref{fig:kv_coverage_map}) odhaluje kritický analytický bod: hustota a~pokrytí nekorelují dokonale.
Francie má hustotu pod \SI{10}{\percent}, ale pokrytí \SI{95}{\percent} --- díky erga omnes; AT má hustotu \char`~30~\%, ale pokrytí \SI{95}{\percent} --- díky odvětvovým zákonným smlouvám.
ČR naopak kombinuje nízkou hustotu s~nízkým pokrytím: žádný z~kompenzačních mechanismů nefunguje.
Pokles odborové hustoty z~\SI{70}{\percent} v~roce 1990 na současných \char`~11~\% tedy neznamená jen pokles počtu členů odborů, ale --- v~absenci erga omnes --- pokles ochrany pracovníků.
Reforma systému by měla cílit primárně na pokrytí (erga omnes), nikoliv na obnovu odborové hustoty jako takové: pokrytí je pro redistribuci produktivity rozhodující proměnnou, hustota je prostředek, nikoliv cíl.

Geografická distribuce odborové hustoty v~EU27 ukazuje na klastrování: severské země a~Belgie udržují hustotu nad \SIrange{50}{60}{\percent}, zatímco středo- a~jihovýchodní Evropa (včetně ČR) se pohybuje pod \SIrange{15}{20}{\percent}.
Srovnání s~mapou \ac{KS} pokrytí (obr.~\ref{fig:kv_coverage_map}) odhaluje kritický analytický bod: hustota a~pokrytí nekorelují dokonale.
Francie má hustotu pod \SI{10}{\percent}, ale pokrytí \SI{95}{\percent} --- díky erga omnes; AT má hustotu \char`~30~\%, ale pokrytí \SI{95}{\percent} --- díky odvětvovým zákonným smlouvám.
ČR naopak kombinuje nízkou hustotu s~nízkým pokrytím: žádný z~kompenzačních mechanismů nefunguje.
Pokles odborové hustoty z~\SI{70}{\percent} v~roce 1990 na současných \char`~11~\% tedy neznamená jen pokles počtu členů odborů, ale --- v~absenci erga omnes --- pokles ochrany pracovníků.
Reforma systému by měla cílit primárně na pokrytí (erga omnes), nikoliv na obnovu odborové hustoty jako takové: pokrytí je pro redistribuci produktivity rozhodující proměnnou, hustota je prostředek, nikoliv cíl.

Geografická distribuce odborové hustoty v~EU27 ukazuje na klastrování: severské země a~Belgie udržují hustotu nad \SIrange{50}{60}{\percent}, zatímco středo- a~jihovýchodní Evropa (včetně ČR) se pohybuje pod \SIrange{15}{20}{\percent}.
Srovnání s~mapou \ac{KS} pokrytí (obr.~\ref{fig:kv_coverage_map}) odhaluje kritický analytický bod: hustota a~pokrytí nekorelují dokonale.
Francie má hustotu pod \SI{10}{\percent}, ale pokrytí \SI{95}{\percent} --- díky erga omnes; AT má hustotu \char`~30~\%, ale pokrytí \SI{95}{\percent} --- díky odvětvovým zákonným smlouvám.
ČR naopak kombinuje nízkou hustotu s~nízkým pokrytím: žádný z~kompenzačních mechanismů nefunguje.
Pokles odborové hustoty z~\SI{70}{\percent} v~roce 1990 na současných \char`~11~\% tedy neznamená jen pokles počtu členů odborů, ale --- v~absenci erga omnes --- pokles ochrany pracovníků.
Reforma systému by měla cílit primárně na pokrytí (erga omnes), nikoliv na obnovu odborové hustoty jako takové: pokrytí je pro redistribuci produktivity rozhodující proměnnou, hustota je prostředek, nikoliv cíl.

Geografická distribuce odborové hustoty v~EU27 ukazuje na klastrování: severské země a~Belgie udržují hustotu nad \SIrange{50}{60}{\percent}, zatímco středo- a~jihovýchodní Evropa (včetně ČR) se pohybuje pod \SIrange{15}{20}{\percent}.
Srovnání s~mapou \ac{KS} pokrytí (obr.~\ref{fig:kv_coverage_map}) odhaluje kritický analytický bod: hustota a~pokrytí nekorelují dokonale.
Francie má hustotu pod \SI{10}{\percent}, ale pokrytí \SI{95}{\percent} --- díky erga omnes; AT má hustotu \char`~30~\%, ale pokrytí \SI{95}{\percent} --- díky odvětvovým zákonným smlouvám.
ČR naopak kombinuje nízkou hustotu s~nízkým pokrytím: žádný z~kompenzačních mechanismů nefunguje.
Pokles odborové hustoty z~\SI{70}{\percent} v~roce 1990 na současných \char`~11~\% tedy neznamená jen pokles počtu členů odborů, ale --- v~absenci erga omnes --- pokles ochrany pracovníků.
Reforma systému by měla cílit primárně na pokrytí (erga omnes), nikoliv na obnovu odborové hustoty jako takové: pokrytí je pro redistribuci produktivity rozhodující proměnnou, hustota je prostředek, nikoliv cíl.

\input{texparts/python/union_density_map}



